\chapter{Upstream's RVV Instruction Set}\label{r04:ch:upstream-rvv}

\ps{}~1.4 contains an RVV instruction set, \isa{rvv}, written by Michael Kuron
(\file{backends/riscv_instruction_sets.py}, helper header
\file{include/riscv_v_helpers.h}; the ratified \code{__riscv_} intrinsic names since
July~2025). It is untouched by this project. This chapter describes what it generates,
because it is both the baseline of every measurement and the starting point of
\isa{rvv_ext}.

\section{The dictionary}

\code{get_vector_instruction_set_riscv(data_type, 'rvv')} builds the dictionary from a
short table of intrinsic stems (\code{'+': 'fadd_vv[0, 1]'}, \code{'loadU':
'le32_v[0]'}, \dots), adding the \code{__riscv_v} prefix, the type suffix
(\code{_f32m1} or \code{_f64m1}) and the active length as last argument. Everything in
it follows from four choices:
\begin{description}[style=nextline]
  \item[Vector-length agnostic] The width is not a number but the C expression
    \code{__riscv_vsetvlmax_e32m1()}: the number of lanes of the machine the code runs on.
  \item[Active length ``elements left''] Every intrinsic gets
    \code{vl = {loop_stop} - {loop_counter}}. Because that placeholder is in the store
    template, the vectorizer keeps the loop whole (\cref{r04:sec:loop}): there is no
    scalar remainder loop, the last iteration simply processes fewer elements.
  \item[LMUL~1] Types \code{vfloat32m1_t}, \code{vfloat64m1_t}, masks \code{vbool32_t}
    or \code{vbool64_t}: one register per vector.
  \item[Plain intrinsics] One intrinsic per operation: \code{vfadd_vv}, \code{vfsub_vv},
    \code{vfmul_vv}, \code{vfdiv_vv}, \code{vfsqrt_v}; unit-stride (\code{vle}/\code{vse})
    and strided (\code{vlse}/\code{vsse}) loads and stores, masked stores; comparisons to
    masks and \code{vmerge} for conditionals; \code{vcpop} for ``any'' and ``all''.
    Scalars are broadcast with \code{vfmv_v_f}.
\end{description}
It also defines \code{'cachelineSize'} and \code{'cachelineZero'} for nontemporal stores
(\cref{r04:sec:rvv-zicboz}).

\section{What it generates}

For the assignment \code{y = 0.5*x + z} on float fields, \isa{rvv} produces this
kernel (unchanged \ps{} output, line breaks added):
\rfourneedspace{9\baselineskip}
\begin{ccode}
FUNC_PREFIX void kernel(float * RESTRICT const _data_x, float * RESTRICT _data_y,
                        float * RESTRICT const _data_z, int64_t const _size_x_0)
{
   for (int64_t ctr_0 = 0; ctr_0 < _size_x_0; ctr_0 += __riscv_vsetvlmax_e32m1())
   {
      __riscv_vse32_v_f32m1(&_data_y[ctr_0],
         __riscv_vfadd_vv_f32m1(
            __riscv_vfmul_vv_f32m1(__riscv_vfmv_v_f_f32m1(0.5f, _size_x_0 - ctr_0),
                                   __riscv_vle32_v_f32m1(& _data_x[ctr_0], _size_x_0 - ctr_0),
                                   _size_x_0 - ctr_0),
            __riscv_vle32_v_f32m1(& _data_z[ctr_0], _size_x_0 - ctr_0),
            _size_x_0 - ctr_0),
         _size_x_0 - ctr_0);
   }
}
\end{ccode}
Three things to notice, because \isa{rvv_ext} changes each of them:
\begin{itemize}
  \item The constant 0.5 is broadcast into a vector (\code{vfmv_v_f}) inside the loop,
    and multiplied vector by vector.
  \item The product and the sum are two intrinsics; RVV intrinsics are opaque to the
    compiler, so it cannot fuse them into one multiply-add.
  \item Every intrinsic carries \code{_size_x_0 - ctr_0}, so the compiler emits
    \code{vsetvli} with that value; nothing is constant.
\end{itemize}

\section{Properties}

\begin{center}\small
\begin{tabularx}{\linewidth}{@{}>{\raggedright\arraybackslash}p{4.2cm}>{\raggedright\arraybackslash}X@{}}
  \toprule
  \textbf{property} & \textbf{consequence}\\
  \midrule
  one binary for every VLEN & checked under QEMU at VLEN 128, 256 and 512 with one
    \code{-march=rv64gcv} binary (\cref{r04:sec:qemu})\\
  no remainder loop & no scalar code in the kernel, and no alignment assumptions\\
  broadcast per use & 476 \code{vfmv_v_f} in the loop body of waLBerla-wind's kernel;
    a broadcast takes 4 cycles on the X60 (measured in \file{docs/notes3})\\
  no FMA & 761 \code{vfmul} and 614 \code{vfadd} in the same loop body\\
  no prefetch, no fixed VLEN, no LMUL above 1 & see \cref{r04:ch:rvv-ext}\\
  no \code{min}/\code{max} templates & no obstacle in practice: the QR model's maximum is
    printed as a compare and a merge (\code{vmfgt}, \code{vmerge})\\
  \bottomrule
\end{tabularx}
\end{center}
The counts are from waLBerla-wind's lattice Boltzmann kernel (D3Q27 cumulant with
Galilean and fourth-order correction, Guo forcing, QR turbulence model, float32, pull),
which \isa{rvv} vectorizes completely in about one minute of code generation: 92 loads,
33 stores and 12 divisions besides the products and sums above.

\section{The active length and the RVV specification}\label{r04:sec:avl}

The active length argument of an intrinsic is an \term{application vector length} (AVL),
not a promise of how many elements are processed. The RVV specification~\cite{rvv-spec}
(section ``Constraints on Setting vl'') lets the hardware choose: for
$\text{AVL} \leq \text{VLMAX}$ it processes AVL elements, for
$\text{AVL} \geq 2\,\text{VLMAX}$ it processes VLMAX, but for
$\text{VLMAX} < \text{AVL} < 2\,\text{VLMAX}$ any value between
$\lceil\text{AVL}/2\rceil$ and VLMAX is allowed. Upstream's loop always advances by
VLMAX. On hardware that took the freedom and split the last two vectors evenly, it
would skip elements. QEMU returns VLMAX in that range, and the QEMU checks use an odd
row length (13 cells), so they cover it. The K1 runs so far use rows of 64 cells, a
multiple of the lane count, so they never reach that range: whether the K1 returns
VLMAX there is not tested. The vector-length specific option of \isa{rvv_ext} does not
depend on it (\cref{r04:sec:ext-vls}).

\section{Nontemporal stores: Zicboz}\label{r04:sec:rvv-zicboz}

With \code{nontemporal}, which waLBerla-wind sets, every unit-stride store gets the
cache-line zeroing guard of \cref{r04:sec:nontemporal}. On RVV the zeroing is the
instruction \code{cbo.zero} of the Zicboz extension~\cite{riscv-cmo}:
\begin{ccode}
inline void cachelineZero(void * p) {        // include/riscv_v_helpers.h
#ifdef __riscv_zicboz
	__asm__ volatile("cbo.zero (%0)"::"r"(p):"memory");
#endif
}
inline size_t cachelineSize() {
#ifdef __riscv_zicboz
	static size_t size = _cachelineSize();   // measured once by zeroing a line
	return size;
#else
	return SIZE_MAX;
#endif
}
\end{ccode}
Both are compiled for Zicboz only if \code{-march} names it. Without it,
\code{cachelineSize()} returns \code{SIZE_MAX}, the alignment test of the guard is never
true, and the guard is dead code: one compare and branch per store. That is the state
of all K1 builds so far, which use \code{-march=rv64gcv_zvl256b}. The K1 does implement
Zicboz (its \code{isa} line in \file{/proc/cpuinfo} lists \code{zicboz}), so building
with \code{-march=rv64gcv_zvl256b_zicboz} is an open experiment.
